\section[Design persistence: the .qrd format]{Design persistence: the \texttt{.qrd} format}\label{sec:serialization}

A rocket design persists as a \texttt{.qrd} file: versioned XML written and read by the two
static entry points of \code{DesignSerializer} (\srcf{core/model/DesignSerializer.h}), on
Boost property tree, the same XML machinery the motor database uses. The format records
exactly what the in-memory model treats as authoritative --- the part tree's typed geometry
parameters, each edge's placement \emph{intent}, the motor's identity, and the
\code{RocketModel}-owned aero options --- and nothing the model derives. No resolved pose,
no composite mass, no absolute coordinate ever reaches disk; the resolver re-derives all of
it on load (\cref{sec:placement}), so a hand-edit that lengthens a tube reloads with every
seam moved to follow the new geometry. Both front-ends reach the format only through these
two calls: the CLI's \code{savedesign}/\code{loaddesign} commands
(\src{cli/Repl.cpp}{1077}, \src{cli/Repl.cpp}{1098}; \cref{sec:cli}) and the GUI's file
menu (\src{gui/MainWindow.cpp}{106}).

The reader is built around one structural commitment: the entire tree is constructed
\emph{detached} --- outside any \code{RocketModel} --- and handed over in a single
\code{installDesign} call at the end. Every validation failure therefore lands before the
current rocket is touched. This section walks the file anatomy, the parameter and link
vocabularies it shares with the factory and the CLI, the motor-by-reference element, and
the load pipeline that turns a file into a live design.

\subsection{File anatomy and the version gate}\label{sec:serialization-anatomy}

The root element is \code{QtRocketDesign} with a \code{version} attribute; the writer
emits \texttt{0.2} (\src{core/model/DesignSerializer.cpp}{198}). The version string is
\emph{major.minor}, and the load gate keys on the major component alone: it splits at the
first dot and rejects any major other than \texttt{0}
(\src{core/model/DesignSerializer.cpp}{233}), so the reader is forward-tolerant across
minor revisions --- a minor bump signals an element-shape change the reader can detect
structurally, not a compatibility break. Under the root sit four elements: a
\code{<design name=...>} attribute element carrying the design's human-facing name, one
recursive \code{<part>} tree, an optional \code{<motor>} reference, and a \code{<sim>}
options block. \Cref{lst:serialization-qrd} shows a complete file from the test corpus.

Each \code{<part>} element carries \code{type} (the part's \code{typeName()}, the exact
factory key --- \texttt{NoseCone}, \texttt{BodyTube}, \texttt{FinSet},
\texttt{HollowSphere}) and \code{name} attributes, a \code{<params>} attribute element, an
optional \code{<link>}, and a \code{<children>} element holding the child \code{<part>}s
in attachment order. The ownership tree in the XML \emph{is} the tree structure; parts are
never cross-referenced by name (names are not unique --- identity is the per-instance id,
which is process-local and deliberately not serialized).

\begin{listing}[htbp]
\begin{minted}{xml}
<?xml version="1.0" encoding="utf-8"?>
<QtRocketDesign version="0.2">
  <design name=""/>
  <part type="NoseCone" name="Nose">
    <params baseRadius="0.0155" length="0.14999999999999999" wallThickness="0"
            density="2500" solid="true"/>
    <children>
      <part type="BodyTube" name="Airframe">   <!-- no <link>: default abut -->
        <params innerRadius="0.0146" outerRadius="0.0155"
                length="0.40000000000000002" density="2500"/>
        <children>
          <part type="FinSet" name="Fins">
            <params density="2500" rootChord="0.059999999999999998"
                    tipChord="0.029999999999999999" span="0.040000000000000001"
                    sweep="0.02" thickness="0.0025000000000000001"
                    bodyRadius="0.0155" finCount="6"/>
            <link seat="OnSurface" parentStation="0" childStation="0" gap="0"/>
            <children/>
          </part>
          <part type="BodyTube" name="Coupler">
            <params innerRadius="0.012" outerRadius="0.0146"
                    length="0.050000000000000003" density="2500"/>
            <link seat="NestInBore" parentStation="1" childStation="0"
                  gap="0.050000000000000003"/>
            <children/>
          </part>
        </children>
      </part>
    </children>
  </part>
  <motor commonName="G80T"/>
  <sim dragCoefficient="1" referenceArea="0.001134"
       referenceAreaOverridden="false"/>
</QtRocketDesign>
\end{minted}
\caption{A complete \texttt{.qrd} file (\srcfcap{tests/data/designs/mid29\_multi.qrd};
attribute lists re-wrapped to the page and the default-abut annotation added, content
otherwise verbatim). All three seat kinds
appear: the airframe abuts the nose by the elided default, the fins seat
\seat{OnSurface}, the coupler nests \seat{NestInBore}. The motor is a database
reference, not an inlined thrust curve.}
\label{lst:serialization-qrd}
\end{listing}

The unpretty decimals are deliberate: the property-tree double translator emits the full
17 significant digits a \cppinline|double| round-trip requires, so a save/load cycle
reproduces geometry --- and therefore mass, CM, and inertia --- bit-exactly. The format
optimizes for faithfulness over prettiness; a human authoring a file by hand writes
\texttt{0.4} and loses nothing.

\subsection{Parameters: the factory round-trip}\label{sec:serialization-params}

The \code{<params>} attributes are not a per-type schema. They are the one shared
vocabulary of \code{PartParams} (\srcf{core/model/parts/PartFactory.h}): every attribute
name in the file is a field of that struct, the writer produces it by calling the
reflector \code{params(part)} (\src{core/model/parts/PartFactory.h}{69}), and the reader
feeds it back to \code{makePart(type, params)}
(\src{core/model/parts/PartFactory.h}{60}), whose documented contract is that
\code{params(makePart(p))} reproduces the geometry. Because the CLI's \code{addpart}
key=value parser fills the same struct, the on-disk attribute names cannot drift from the
interactive vocabulary --- there is exactly one spelling of \code{outerRadius} in the
system, and it is a struct field, not a string constant in two parsers.

Every field is \cppinline|std::optional|, and \code{writeParams} emits only the fields
that are present (\src{core/model/DesignSerializer.cpp}{32}), so a nose cone's element
carries cone attributes and nothing else. On read, a missing key stays
\cppinline|std::nullopt| --- distinct from a supplied zero --- and \code{makePart} applies
its documented defaults or throws for a missing required field, the same fail-closed path
an interactive command takes. The one non-numeric field, \code{solid}, serializes as the
strings \texttt{true}/\texttt{false}, matching the motor database's XML idiom.

\subsection[Placement: the <link> element, elided by default]{Placement: the \texttt{<link>} element, elided by default}
\label{sec:serialization-link}

A child's spatial relationship to its parent is its \cppinline|StationLink|
(\cref{sec:placement}), serialized as an inline \code{<link>} on the child with exactly
the four authored quantities: \code{seat}, \code{parentStation}, \code{childStation},
\code{gap}. The seat enumerator maps to its attribute string through the single shared
table \code{seatKindToString}/\code{seatKindFromString}
(\src{core/model/parts/Placement.h}{41}); the writer
(\src{core/model/DesignSerializer.cpp}{57}), the reader
(\src{core/model/DesignSerializer.cpp}{140}), and the CLI's \code{seat=} token parser
(\src{cli/Repl.cpp}{232}) all consult it, so the three agree on the spellings
(\seat{Abut}, \seat{NestInBore}, \seat{OnSurface}) by construction rather than by three
independently maintained switch statements. \code{seatKindFromString} returns
\cppinline|std::nullopt| for an unknown name, and the reader converts that into a thrown
error on the same path as \code{makePart}'s unknown-type rejection --- a hand-edited seat
typo is a refused file with a directed diagnostic, never a silent coercion to \seat{Abut}.

The link's quaternion (\code{childRot}) is not persisted: it is identity in 3-DOF, so a
stored link is fully described by its scalars and seat, and the file format acquires the
6-DOF field only when a non-identity value exists to store.

\begin{designrule}[The common case costs zero authored numbers]
\code{writePart} compares each link against a default-constructed
\cppinline|StationLink| (\code{isDefaultLink},
\src{core/model/DesignSerializer.cpp}{67}) and elides the element entirely when they are
equal (\src{core/model/DesignSerializer.cpp}{85}); the reader, seeing no \code{<link>},
attaches by the same default (\src{core/model/DesignSerializer.cpp}{175}). The default is
a meaningful physical relationship --- child fore plane abutting parent aft plane, the
natural nose$\to$body$\to\ldots$ stack --- so eliding it loses nothing, and the two halves
of the contract (elide on write, default on read) are pinned to the same
default-constructed struct. The payoff is visible in \cref{lst:serialization-qrd}: the
airframe, the largest part in the design, is attached by zero authored placement numbers,
and the file's size tracks the design's \emph{placement complexity}, not its part count.
\end{designrule}

Within major version \texttt{0}, one element shape is refused by name: a \texttt{0.1}-era
file records placement as a resolved CM-to-CM \code{<offset>} coordinate triple, which
the intent-only model cannot faithfully represent. The reader throws on that element with
a message directing the user to re-create the design
(\src{core/model/DesignSerializer.cpp}{180}) --- mis-seating every child by silently
treating a resolved coordinate as the abut default would corrupt the design while
appearing to load it.

\subsection{The motor: a database reference}\label{sec:serialization-motor}

The motor element is one attribute --- \code{<motor commonName="G80T"/>}
(\src{core/model/DesignSerializer.cpp}{205}) --- plus an optional \code{<link>} for its
seat, elided under the same default-equal rule as any other edge
(\src{core/model/DesignSerializer.cpp}{207}). The file stores a \emph{reference} into the
motor database, never an inlined thrust curve.

\begin{designrule}[The database is the source of truth for motors]
A motor is certified vendor data --- a thrust curve, propellant and total masses, burn
time --- that the design merely selects (\cref{sec:motor}). Inlining a copy into every
\texttt{.qrd} would create as many divergent copies of that data as there are saved
designs, each frozen at its save date; resolving by common name on every load means a
corrected or updated database benefits every existing file. It also keeps the file small
and portable: a \texttt{.qrd} carries authored intent measured in tens of numbers, not a
hundred-sample thrust curve.
\end{designrule}

The save walk enforces the reference form structurally: \code{writePart} skips any
\code{Motor} child while recursing (\src{core/model/DesignSerializer.cpp}{95}), so the
motor --- an ordinary node in the live tree (\cref{sec:tree}) --- never appears as a
\code{<part>}, only as the \code{<motor>} element. The reader enforces the inverse:
\code{buildNode} rejects \code{type="Motor"} outright
(\src{core/model/DesignSerializer.cpp}{155}), because the geometry factory cannot rebuild
a motor from \code{<params>} --- a \code{Motor} tree part can only occur in a hand-edited
or foreign file, and refusing it names the correct channel in the diagnostic.

On load, the common name is resolved against the caller-supplied
\code{MotorModelDatabase} (\src{core/model/DesignSerializer.cpp}{246}). A hit attaches a
fresh \code{Motor} node under the root with the persisted (or defaulted) link
(\src{core/model/DesignSerializer.cpp}{253}); a miss logs a warning and the design loads
without a motor (\src{core/model/DesignSerializer.cpp}{257}) --- degraded, flyable after
the user selects a motor, and fully intact in geometry.

\begin{rejected}[Rejected --- fail-closed motor resolution]
Symmetry with the seat and type gates argues for throwing on a motor miss; the reference
form argues the other way, and wins. A seat typo is evidence the file is corrupt, but a
missing motor is evidence about the \emph{machine}: the same valid file loaded on a
workstation whose database lacks one vendor file would be refused in its entirety ---
geometry, placements, and all --- over the one datum the database can supply later.
Fail-closed here would make every \texttt{.qrd} silently non-portable across database
states. The warning path degrades exactly as far as the missing information requires, and
\cref{sec:flight}'s launch gate (\code{isMotorSet()}) already prevents flying the
degraded design by accident.
\end{rejected}

\subsection{Load: build detached, install once}\label{sec:serialization-load}

\Cref{fig:serialization-load} traces the load pipeline;
\cref{lst:serialization-loadtail} is its tail in code. After
\cppinline|read_xml| and the major-version gate, \code{buildNode} recursion constructs
the entire tree as detached \code{PartNode}s --- \code{PartNode::make} wrapping each
factory-built part, \code{addChild} assembling bottom-up
(\src{core/model/DesignSerializer.cpp}{240}), the same detached-building path a clone
uses (\cref{sec:tree}). The motor, when resolved, is attached to that detached root the
same way. Only then does the one mutation of shared state happen:
\code{rocket.installDesign(newRoot)} (\src{core/model/DesignSerializer.cpp}{262}), which
atomically replaces the tree through \code{PartsModel::installRoot} and re-resolves the
motor borrow inside the install facade.

\tikzfig{activity-load}{The load pipeline. Every fail-closed exit --- version, part type,
seat kind, placement element --- throws from inside the detached region, before
\code{installDesign}; the motor miss is the one soft path. The sim attributes apply after
the install so a persisted manual reference area survives the install-time override
reset.}{fig:serialization-load}

\begin{keyinsight}[A malformed file leaves the current rocket untouched]
Nothing before \code{installDesign} touches the \code{RocketModel}. Every validation
failure --- unreadable XML, unknown major version, unknown part type, non-physical
geometry thrown by a concrete constructor, unknown seat, a \code{Motor} tree part, a
\texttt{0.1} \code{<offset>} --- propagates out of the detached build and destroys only
the half-built detached tree. A user who mistypes a filename or opens a corrupt file gets
an error message and keeps the design they were editing; there is no partially-loaded
state to recover from, by construction rather than by rollback code.
\end{keyinsight}

\begin{listing}[htbp]
\begin{minted}{cpp}
// Build the whole tree (detached) before touching the rocket, then install in one
// shot so a malformed file leaves the existing rocket untouched.
std::unique_ptr<PartNode> newRoot = buildNode(root.get_child("part"));

// Motor by common name, re-resolved against the supplied database (a miss warns,
// never throws). An ordinary node: attached under the root with its persisted link.
if(const auto name = root.get_optional<std::string>("motor.<xmlattr>.commonName"))
{
    if(const auto m = motors.getMotorModel(*name))
    {
        part::StationLink motorLink{};
        if(root.get_child_optional("motor.link"))
        {
            motorLink = readLink(root.get_child("motor"));
        }
        newRoot->addChild(PartNode::make(std::make_unique<part::Motor>("Motor", *m)),
                          motorLink);
    }
    else
    {
        utils::Logger::getInstance()->warn(/* ... loaded without a motor */);
    }
}

rocket.installDesign(std::move(newRoot));
rocket.setName(root.get<std::string>("design.<xmlattr>.name", ""));

// Apply the sim/aero block after the install so a restored manual override wins
// over its reset.
rocket.setDragCoefficient(
    root.get<double>("sim.<xmlattr>.dragCoefficient", rocket.getDragCoefficient()));
const bool overridden =
    root.get<std::string>("sim.<xmlattr>.referenceAreaOverridden", "false") == "true";
if(overridden)
{
    rocket.setReferenceArea(
        root.get<double>("sim.<xmlattr>.referenceArea", rocket.getReferenceArea()));
}
\end{minted}
\caption{The tail of \code{DesignSerializer::load}
(\srccap{core/model/DesignSerializer.cpp}{240}, lightly elided): detached motor attach,
the single atomic install, then name and sim attributes.}
\label{lst:serialization-loadtail}
\end{listing}

\subsection[The <sim> block and the override tri-state]{The \texttt{<sim>} block and the override tri-state}
\label{sec:serialization-sim}

The \code{<sim>} element persists only the \code{RocketModel}-owned aero options:
\code{dragCoefficient}, \code{referenceArea}, and \code{referenceAreaOverridden}
(\src{core/model/DesignSerializer.cpp}{216}). Launch and environment options ---
rail velocity, atmosphere, gravity, integrator --- are session configuration owned by the
front-ends (\cref{sec:cli}), not rocket state, and deliberately stay out of the design
file: a design describes a vehicle, not a flight.

The reference area needs two attributes because it is one value with two provenances
(\cref{sec:flight}): a geometry-derived default (the widest frontal disc in the tree) or
a manual override set through \code{setReferenceArea}, which latches the
\code{referenceAreaOverridden} flag (\src{core/model/RocketModel.h}{69}). Persisting the
flag alongside the value gives the reader three distinct behaviors: with
\code{referenceAreaOverridden="true"} it calls \code{setReferenceArea}, restoring both
the value and the manual-override provenance; with \texttt{"false"} it applies nothing
--- the stored \code{referenceArea} is a record of the value at save time, and the
rocket's existing area continues to govern; with the attribute absent it likewise leaves the
rocket's current state alone. Storing the value unconditionally but applying it
conditionally keeps the file self-describing while never promoting a derived number into
authored intent.

Ordering carries the correctness burden: \code{installDesign} clears the override flag as
an install side effect --- a freshly-installed airframe must not inherit a manual area
sized for a different one (\src{core/model/RocketModel.cpp}{121}) --- so the loader
applies the sim block \emph{after} the install (\src{core/model/DesignSerializer.cpp}{266}).
A persisted override is thereby re-established on top of the reset rather than clobbered
by it; applying the attributes first would make every load silently discard the user's
saved area. The CLI then syncs its staged session copies from the freshly loaded rocket
(\src{cli/Repl.cpp}{1107}), so the next \code{fly} sees the restored values.
